\newpage

% ######################################################################
% PART XV: EXERCISES & WHAT-IF
% ######################################################################
\part{Exercises \& What-If Scenarios}

% ######################################################################
% SECTION 59: app.py + Models Full Walkthrough
% ######################################################################
\section{app.py + Models --- Jo Abhi Tak Full Nahi Dekha}

\subsection{app.py --- Session State Flow}

\begin{codestorybox}
\begin{lstlisting}[style=pythonstyle, caption={app.py --- key flow}]
# 1. Page config + theme
st.set_page_config(page_title="Multi-Agent Research", layout="wide")

# 2. Session state lazy init
if "theme" not in st.session_state:
    st.session_state.theme = "light"
if "messages" not in st.session_state:
    st.session_state.messages = []
if "final_report" not in st.session_state:
    st.session_state.final_report = ""

# 3. Sidebar: theme picker + PDF upload
theme = st.sidebar.selectbox("Theme", ["light", "dark"],
                             index=["light","dark"].index(st.session_state.theme))
if theme != st.session_state.theme:
    st.session_state.theme = theme
    st.rerun()          # theme apply karne ke liye rerun

uploaded = st.sidebar.file_uploader("Upload PDF", type=["pdf"])
if uploaded:
    save + ingest -> st.success(f"{n} chunks added")

# 4. Topic input + Run
topic = st.text_input("Research topic")
if st.button("Run") and topic:
    # reset state
    st.session_state.messages = []
    # stream pipeline
    for event in research_graph.stream(initial_state(topic)):
        for node, output in event.items():
            st.session_state.agent_statuses[node] = "done"
            with st.chat_message(node):
                st.write(output.get("messages", [""])[-1])
    st.session_state.final_report = ... # last state se

# 5. Report render + download
if st.session_state.final_report:
    st.markdown(st.session_state.final_report)
    st.download_button("Download MD", ..., file_name="report.md")
\end{lstlisting}

\textbf{Interview points:} (1) \texttt{st.rerun()} theme change par --- Streamlit idiom; (2) \texttt{file\_uploader} + ingest --- RAG loop close; (3) chat-message styling se agent progress visual; (4) \texttt{download\_button} --- result export.
\end{codestorybox}

\subsection{backend/api/models.py --- Full Walkthrough}

\begin{codestorybox}
\begin{lstlisting}[style=pythonstyle, caption={models.py --- dono models}]
from django.db import models
from django.contrib.auth.models import User

class ResearchSession(models.Model):
    topic = models.CharField(max_length=500)
    report = models.TextField()
    messages = models.JSONField(default=list)     # agent log
    user = models.ForeignKey(User, null=True, blank=True,
                             on_delete=models.SET_NULL)
    sources_count = models.IntegerField(default=0)
    revision_count = models.IntegerField(default=0)
    created_at = models.DateTimeField(auto_now_add=True)

    def __str__(self):
        return f"Session: {self.topic[:50]}"

class BenchmarkEvaluation(models.Model):
    topic = models.CharField(max_length=500)
    single_agent_report = models.TextField()
    multi_agent_report = models.TextField()
    single_agent_depth = models.IntegerField()
    multi_agent_depth = models.IntegerField()
    single_agent_verifiability = models.IntegerField()
    multi_agent_verifiability = models.IntegerField()
    verdict = models.CharField(max_length=40)
    created_at = models.DateTimeField(auto_now_add=True)

    def __str__(self):
        return f"Benchmark: {self.topic[:50]} ({self.verdict})"
\end{lstlisting}

\textbf{Design decisions:}
\begin{itemize}
\item \texttt{JSONField} messages --- flexible list storage (no separate table needed).
\item \texttt{on\_delete=models.SET\_NULL} --- user delete ho to session rahe (data loss nahi).
\item \texttt{null=True, blank=True} --- anonymous sessions allowed.
\item \texttt{auto\_now\_add} --- created\_at automatic, immutable.
\item IntegerFields for scores --- 0-10 scale, DB-level simplicity.
\end{itemize}
\end{codestorybox}

% ######################################################################
% SECTION 60: BUILD-IT-YOURSELF EXERCISES
% ######################################################################
\section{Build-It-Yourself --- Mini Versions With Solutions}

\begin{storybox}
Interview practice ke liye: project ke 4 core components ko \textbf{zero-dependency} mini versions mein implement karo (30-45 min each). Ye exercises aapki understanding ko prove karte hain --- agar mini version bana sakte ho, to bada version samajhte ho.
\end{storybox}

\subsection{Exercise 1: Mini Vector Store}

\begin{codestorybox}
\textbf{Task:} Bina ChromaDB ke, ek class jo documents store kare aur cosine similarity se top-k retrieve kare. \textbf{Hint:} numpy dot products. \textbf{Solution:}

\begin{lstlisting}[style=pythonstyle, caption={mini vector store}]
import numpy as np

class MiniVectorStore:
    def __init__(self):
        self.docs = []      # list[str]
        self.vectors = []   # list[np.ndarray]

    def add(self, text: str, vector: np.ndarray):
        self.docs.append(text)
        self.vectors.append(vector / np.linalg.norm(vector))

    def search(self, query_vec: np.ndarray, k: int = 3) -> list[str]:
        q = query_vec / np.linalg.norm(query_vec)
        scores = [float(v @ q) for v in self.vectors]  # cosine
        idx = sorted(range(len(scores)), key=lambda i: -scores[i])[:k]
        return [self.docs[i] for i in idx]

# use:
# store.add("RAG is retrieval augmented generation", embed("RAG..."))
# store.search(embed("what is RAG"), k=2)
\end{lstlisting}

\textbf{Key concepts proven:} cosine similarity = normalized dot product; top-k sorting; O(n) brute force (HNSW production mein).
\end{codestorybox}

\subsection{Exercise 2: Mini Agent Loop}

\begin{codestorybox}
\textbf{Task:} Bina LangGraph ke, ek simple loop: research $\rightarrow$ analyze $\rightarrow$ verify $\rightarrow$ revise (max 3). \textbf{Solution:}

\begin{lstlisting}[style=pythonstyle, caption={mini agent loop}]
def run_pipeline(topic, call_llm, search):
    state = {"topic": topic, "sources": [], "analysis": "",
             "passed": False, "revisions": 0}
    while not state["passed"] and state["revisions"] <= 3:
        state["sources"] = search(topic)          # researcher
        state["analysis"] = call_llm("analyze", state["sources"])
        verdict = call_llm("verify", state["analysis"], state["sources"])
        state["passed"] = "PASSED" in verdict.upper()
        state["revisions"] += 1
    return {"report": call_llm("write", state["analysis"], state["sources"]),
            "revisions": state["revisions"]}
\end{lstlisting}

\textbf{Interview question:} ``LangGraph isse kya better karta hai?'' --- streaming, visualization, parallelism, checkpointing, declarative flow.
\end{codestorybox}

\subsection{Exercise 3: Mini Benchmark Judge}

\begin{codestorybox}
\textbf{Task:} Do reports ko score karne wala heuristic judge (bina LLM ke). \textbf{Solution:}

\begin{lstlisting}[style=pythonstyle, caption={mini heuristic judge}]
import re

def judge(report: str) -> dict:
    sections = len(re.findall(r'^#{1,3} ', report, re.M))
    bullets = len(re.findall(r'^[-*] ', report, re.M))
    words = len(report.split())
    citations = len(re.findall(r'\[\d+\]|https?://', report))
    depth = min(10, sections // 3 + bullets // 5 + words // 300)
    verifiability = min(10, citations * 2)
    return {"depth": depth, "verifiability": verifiability}
\end{lstlisting}

\textbf{Interview point:} Heuristics deterministic hain (tests ke liye), LLM judge quality deta hai --- is project ka evaluator dono combine karta hai (LLM first, heuristic fallback).
\end{codestorybox}

% ######################################################################
% SECTION 61: WHAT-IF SCENARIOS
% ######################################################################
\section{What-If Scenarios --- Aage Badhne Ke Raaste}

\subsection{Scenario 1: Agent Memory Add Karna}

\begin{conceptbox}
\textbf{Problem:} Agents har run mein past se forget karte hain (sirf rag\_context\_past\_research milta hai).

\textbf{Solution sketch:} (1) \textbf{Short-term:} LangGraph checkpointing --- har run ki state save, resume; (2) \textbf{Long-term:} Agent-specific summaries ChromaDB mein (``what did analyst learn about topic X''), retrieval on next run; (3) \textbf{Conversation memory:} add\_messages history pipeline ko. Interview answer: ``Memory = retrieval over past states + summaries, ek flag ke saath enable/disable.''
\end{conceptbox}

\subsection{Scenario 2: Planner Agent Add Karna}

\begin{conceptbox}
\textbf{Problem:} Complex topics (``AI in healthcare AND finance'') ek hi researcher pass mein poori nahi cover hote.

\textbf{Solution sketch:} (1) \textbf{Planner node} graph ke START par --- topic ko sub-tasks mein toda (``healthcare AI'', ``finance AI'', ``joint analysis''); (2) Har sub-task ka apna mini-pipeline (parallel fan-out); (3) \textbf{Aggregator} node --- sub-reports ko final report mein merge. Interview point: ``Hierarchical multi-agent --- plan $\rightarrow$ execute $\rightarrow$ aggregate. LangGraph fan-out/fan-in edges se possible.''
\end{conceptbox}

\subsection{Scenario 3: Multi-Hop RAG}

\begin{conceptbox}
\textbf{Problem:} Answer ke liye 2+ documents ka combination chahiye (``Company X ka revenue kaun sa product se aaya'' --- pehle product list, phir revenue).

\textbf{Solution sketch:} (1) Round 1: retrieve + answer attempt with confidence; (2) Low confidence $\rightarrow$ extracted entities se follow-up query; (3) Round 2: retrieve + combine. Loop guard max 3. Interview line: ``Naive RAG single-hop hai; multi-hop query decomposition se complex reasoning improve hoti hai --- same revision-loop pattern se implement hota hai jo fact-checker mein hai.''
\end{conceptbox}

\subsection{Scenario 4: Real-Time Streaming Reports}

\begin{conceptbox}
\textbf{Problem:} 176s wait --- user ko real-time partial report dikhao.

\textbf{Solution sketch:} (1) Har node ka output stream (already \texttt{graph.stream}); (2) UI mein node-by-node sections render (researcher ke sources, analyst brief, etc.); (3) Writer ka \textbf{token-level streaming} (Groq streaming) --- last section live type hote dikhe. Interview point: ``Progressive disclosure --- user ko har step ka value dikhta hai, wait boring nahi.''
\end{conceptbox}

\subsection{Scenario 5: Multi-Language Support}

\begin{conceptbox}
\textbf{Problem:} Hindi/Hinglish queries par reports English mein.

\textbf{Solution sketch:} (1) \textbf{Embedding model:} multilingual (multilingual-e5 / bge-m3) --- cross-lingual retrieval; (2) \textbf{Language param:} writer prompt mein ``report language = Hindi/English''; (3) \textbf{Translation layer} optional. Interview point: ``Embedding model change = retrieval quality change --- benchmark se verify karna zaroori (Retrieval eval), assumptions nahi.''
\end{conceptbox}

\begin{memorybox}
\textbf{Final exercise challenge:} In 5 scenarios mein se 2 ko implement karke demo banao (even paper sketch enough) --- interview mein ``I've thought about scaling this'' ka proof milega. Koi bhi ``what if'' question aaye --- answer ready hai.
\end{memorybox}
